\Askhsh{28299_SOLUTION}{1}
\begin{ylh}{1}
    \item [Ύλη από εικόνα]
\end{ylh}
\begin{wrapfigure}[20]{r}{0.5\textwidth} % The figure will float to the right
\centering
\begin{tikzpicture}
\begin{axis}[
    axis lines = middle,
    xlabel = {$x$},
    ylabel = {$y$},
    xmin = -2.5, xmax = 4.5,
    ymin = -2.5, ymax = 4.5,
    xtick = {-2,-1,0,1,2,3,4},
    ytick = {-2,-1,0,1,2,3,4},
    grid = both,
    grid style = {gray!30},
    width = 0.48\textwidth, height = 0.48\textwidth,
    tick label style = {font=\scriptsize},
    every axis x label/.style = {at={(current axis.right of origin)}, anchor=north west},
    every axis y label/.style = {at={(current axis.above origin)}, anchor=south},
    clip=false,
]
    % Line y = x (dashed)
    \addplot[dashed, gray, very thick, domain=-2.5:4.5] {x};
    \node[below left, font=\scriptsize] at (axis cs:-0.5,-0.5) {$y=x$};

    % Plot for y = f(x) (dark grey)
    \addplot[very thick, gray] plot[smooth] coordinates {(-1,-1) (0,1) (1,2) (3,3) (4,4)};
    \node[right, gray, font=\scriptsize] at (axis cs:1,2) {$y = f(x)$};

    % Plot for y = f^{-1}(x) (blue/secondarycolor)
    \addplot[very thick, secondarycolor] plot[smooth] coordinates {(-1,-1) (1,0) (2,1) (3,3) (4,4)};
    \node[above, secondarycolor, font=\scriptsize] at (axis cs:2,1) {$y = f^{-1}(x)$};

    % Plot all the points as black dots
    \addplot[only marks, mark=*, mark size=2pt, black] coordinates {
        (-1,-1) (0,1) (1,2) (3,3) (4,4) % Points on f(x)
        (1,0) (2,1)                   % Additional points on f^-1(x) (symmetric to (0,1) and (1,2))
    };
\end{axis}
\end{tikzpicture}
\caption*{Οι γραφικές παραστάσεις $C_f$ και $C_{f^{-1}}$}
\end{wrapfigure}

\OnPar{ΛΥΣΗ}
\begin{alist}
    \item [α)] Παρατηρούμε ότι κάθε οριζόντια ευθεία τέμνει την γραφική παράσταση της $f$ το πολύ σε ένα σημείο. Αυτό σημαίνει ότι δεν υπάρχουν σημεία της γραφικής παράστασης με την ίδια τεταγμένη. Επομένως η $f$ είναι $1-1$, άρα ορίζεται η αντίστροφη συνάρτηση $f^{-1}$ της $f$.

    \item [β)] Σχεδιάζουμε στο σχήμα την ευθεία $y = x$. Τα σημεία τομής της $C_f$ με την ευθεία $y = x$ είναι τα σημεία της $C_f$ των οποίων η τετμημένη είναι ίση με την τεταγμένη άρα είναι τα σημεία $(4,4)$, $(3,3)$ και $(-1,-1)$.

    \item [γ)] Η γραφική παράσταση $C_{f^{-1}}$ της $f^{-1}$ είναι συμμετρική της γραφικής παράστασης $C_f$ της $f$, ως προς την ευθεία $y = x$. Επομένως η $C_{f^{-1}}$ διέρχεται από τα σημεία $(4,4)$, $(3,3)$ και $(-1,-1)$ καθώς και από τα σημεία $(1,0)$ και $(2,1)$ που είναι τα συμμετρικά των σημείων $(0,1)$ και $(1,2)$ της $C_f$ ως προς την ευθεία $y = x$. Οι γραφικές παραστάσεις $C_f$ και $C_{f^{-1}}$ φαίνονται παρακάτω:
\end{alist}